% 02-phuong-phap.tex — Chương 2: FCM, PSO và các phiên bản, chỉ số đánh giá, PSO tìm tâm (hướng ban đầu), trọng số
% thuộc tính trong FCM, đề xuất PSO tìm trọng số thuộc tính. Tham số đúng theo code/*.py; không có số liệu kết quả.
\chapter{FUZZY C-MEANS, PSO VÀ ĐỀ XUẤT PSO TÌM TRỌNG SỐ THUỘC TÍNH}
\label{chap:phuong-phap}

Chương này có 2 phần. Phần nền gồm các thuật toán và thước đo dùng chung cho cả 2 giai đoạn nghiên cứu:
Fuzzy C-Means (FCM) và lý do kết quả của FCM có thể phụ thuộc khởi tạo (mục~\ref{sec:fcm}), thuật toán tối ưu bầy đàn
(Particle Swarm Optimization, PSO) ở dạng gốc của Kennedy và Eberhart~\cite{kennedy1995pso} (mục~\ref{sec:pso-goc}), 4
phiên bản cải tiến lấy từ các bài đã đọc (mục~\ref{sec:pso-cai-tien}) và các chỉ số đánh giá (mục~\ref{sec:chi-so}). Phần
sau đi theo thứ tự nhóm thực hiện: mục~\ref{sec:pso-khoi-tao} là hướng thử đầu tiên, dùng PSO tìm tâm khởi tạo cho FCM; khi
hướng này không cải thiện FCM trên dữ liệu tín dụng, nhóm chuyển sang trọng số thuộc tính. Mục~\ref{sec:trong-so-fcm} chỉ ra
vì sao FCM không tự chọn được trọng số và giới hạn của cách tính trọng số bằng công thức; mục~\ref{sec:de-xuat} là đề xuất cải
tiến của nhóm, dùng PSO tìm trọng số thuộc tính cho FCM.

Các thuật toán trong tiểu luận không độc lập với nhau (Bảng~\ref{tab:ch2-ban-do}). Mọi phương pháp đều là FCM, chỉ khác ở
hai chỗ có thể thay: \emph{tâm khởi tạo} và \emph{trọng số thuộc tính}. PSO là một bộ tìm kiếm duy nhất, lắp vào một trong hai
chỗ đó; 5 phiên bản PSO dùng chung một khung và chỉ khác nhau ở quy tắc di chuyển hạt. Một thí nghiệm vì vậy được xác định
bằng 3 lựa chọn: cách khởi tạo tâm, cách đặt trọng số và, nếu có PSO, phiên bản PSO.

\begin{table}[htbp]
\centering
\caption{Bản đồ các thuật toán: một lõi FCM, hai chỗ thay và một bộ tìm kiếm PSO}
\label{tab:ch2-ban-do}
\footnotesize
\setlength{\tabcolsep}{3pt}
\renewcommand{\arraystretch}{1.2}
\begin{tabular}{@{}>{\raggedright\arraybackslash}p{2.5cm}>{\raggedright\arraybackslash}p{7.0cm}>{\raggedright\arraybackslash}p{1.6cm}>{\raggedright\arraybackslash}p{3.6cm}@{}}
\toprule
Khối & Các lựa chọn & Mục & Tệp mã \\
\midrule
Lõi: FCM & Lặp cập nhật độ thuộc $U$ và tâm $V$ tới hội tụ; mọi phương pháp kết thúc bằng khối này &
\ref{sec:fcm} & \texttt{fcm.py} \\
\addlinespace
Chỗ thay 1: tâm khởi tạo & $U$ ngẫu nhiên; $c$ điểm ngẫu nhiên; k-means++; \textbf{PSO} (hạt là một bộ tâm, fitness $R_m$) &
\ref{sec:pso-khoi-tao} & \texttt{pso\_khoi\_tao.py} \\
\addlinespace
Chỗ thay 2: trọng số thuộc tính & $\symbf{\lambda} = \mathbf{1}$ (FCM thường); WFCM giải tích; tìm ngẫu nhiên; \textbf{PSO} (hạt là
vector $\symbf{\lambda}$, fitness Silhouette cộng phạt; đề xuất của nhóm) & \ref{sec:trong-so-fcm}, \ref{sec:de-xuat} &
\texttt{cach\_ket\_hop.py}, \texttt{pso\_trong\_so.py} \\
\addlinespace
Bộ tìm kiếm PSO & Một vòng lặp chung. PSO gốc, PSO-W, PSO-S, PSO-LOG chỉ khác hàm quán tính $\omega(t)$; QPSO thay cập nhật vận
tốc bằng lấy mẫu quanh điểm hút & \ref{sec:pso-goc}, \ref{sec:pso-cai-tien} & \texttt{pso\_co\_so.py}, \texttt{thuat\_toan\_a.py},
\texttt{thuat\_toan\_b.py} \\
\addlinespace
Đánh giá & Chính: $\Delta r$, AUC phân hoạch; phụ: Silhouette, PC, PE, XB, thời gian; so sánh bằng Wilcoxon kèm Holm &
\ref{sec:chi-so}, \ref{sec:kiem-dinh} & \texttt{danh\_gia.py}, \texttt{chi\_so\_danh\_gia.py}, \texttt{phan\_tich.py} \\
\bottomrule
\end{tabular}
\end{table}

Trong toàn chương, dữ liệu là tập $X = \{\mathbf{x}_1, \dots, \mathbf{x}_n\} \subset \mathbb{R}^D$ đã qua tiền xử lý
(mục~\ref{sec:tien-xu-ly}), $c$ là số cụm, $U = [u_{ik}]$ là ma trận độ thuộc kích thước $n \times c$ và
$V = [\mathbf{v}_1, \dots, \mathbf{v}_c]^{\top}$ là ma trận tâm cụm kích thước $c \times D$. Khoảng cách là khoảng cách
Euclid, $d_{ik} = \lVert \mathbf{x}_i - \mathbf{v}_k \rVert$.

% ------------------------------------------------------------
\section{Thuật toán Fuzzy C-Means}
\label{sec:fcm}

\subsection{Hàm mục tiêu và điều kiện tối ưu}

FCM~\cite{bezdek1981pattern,bezdek1984fcm} cho mỗi điểm thuộc về mọi cụm với mức độ trong $[0, 1]$, tổng bằng một, và tìm
$(U, V)$ cực tiểu hoá
\begin{equation}
  J_m(U, V) = \sum_{i=1}^{n} \sum_{k=1}^{c} u_{ik}^{m} \, d_{ik}^{2},
  \label{eq:fcm-j}
\end{equation}
với ràng buộc
\begin{equation}
  u_{ik} \in [0, 1], \qquad \sum_{k=1}^{c} u_{ik} = 1 \;\; (i = 1, \dots, n), \qquad
  0 < \sum_{i=1}^{n} u_{ik} < n \;\; (k = 1, \dots, c).
  \label{eq:fcm-rang-buoc}
\end{equation}
Số mũ $m > 1$ điều khiển độ mờ: khi $m \to 1$, FCM tiến về k-means cứng; khi $m$ tăng, độ thuộc tiến về $1/c$. Dùng nhân
tử Lagrange cho ràng buộc tổng, điều kiện cần để $(U, V)$ là điểm dừng của~\eqref{eq:fcm-j} là~\cite{bezdek1981pattern}
\begin{align}
  \mathbf{v}_k &= \frac{\sum_{i=1}^{n} u_{ik}^{m} \mathbf{x}_i}{\sum_{i=1}^{n} u_{ik}^{m}}, \label{eq:fcm-v}\\[2pt]
  u_{ik} &= \left[ \sum_{l=1}^{c} \left( \frac{d_{ik}}{d_{il}} \right)^{2/(m-1)} \right]^{-1}. \label{eq:fcm-u}
\end{align}
Khi $d_{ik} = 0$, mã nguồn thay khoảng cách bằng một số dương rất nhỏ để tránh chia cho không.

FCM luân phiên tính $V$ theo~\eqref{eq:fcm-v} và $U$ theo~\eqref{eq:fcm-u} tới khi thay đổi đủ nhỏ
(Thuật toán~\ref{alg:fcm}, \texttt{fcm.py}). Khi khởi động từ một bộ tâm (c điểm dữ liệu, k-means++~\cite{arthur2007kmeanspp} hoặc gbest của PSO),
FCM cập nhật $U$ trước và dừng theo thay đổi của tâm. Mọi cách khởi tạo bằng tâm trong tiểu luận (điểm ngẫu nhiên,
k-means++, gbest của PSO) kết thúc bằng cùng một FCM với cùng tiêu chuẩn dừng, nên chúng chỉ khác nhau ở điểm xuất phát.
Riêng đối chứng FCM khởi tạo $U$ ngẫu nhiên dừng theo thay đổi của $U$, đúng dạng gốc.

\begin{algorithm}[htbp]
\caption{Fuzzy C-Means (bản cài đặt trong \texttt{fcm.py})}
\label{alg:fcm}
\begin{algorithmic}[1]
\REQUIRE Dữ liệu $X$, số cụm $c$, số mũ $m = 2$, ngưỡng $\varepsilon = 10^{-5}$, số vòng tối đa $I_{\max} = 10.000$,
  tâm ban đầu $V^{(0)}$ (tuỳ chọn)
\ENSURE Ma trận độ thuộc $U$, tâm cụm $V$
\IF{không có $V^{(0)}$}
  \STATE Khởi tạo $U^{(0)}$ ngẫu nhiên, chuẩn hoá mỗi hàng về tổng bằng 1
\ENDIF
\FOR{$t = 0, 1, \dots, I_{\max} - 1$}
  \IF{khởi động từ $U$}
    \STATE $V^{(t+1)} \leftarrow$ \eqref{eq:fcm-v} từ $U^{(t)}$; \quad $U^{(t+1)} \leftarrow$ \eqref{eq:fcm-u} từ $V^{(t+1)}$
    \STATE \textbf{nếu} $\max |U^{(t+1)} - U^{(t)}| < \varepsilon$ \textbf{thì} dừng
  \ELSE
    \STATE $U^{(t+1)} \leftarrow$ \eqref{eq:fcm-u} từ $V^{(t)}$; \quad $V^{(t+1)} \leftarrow$ \eqref{eq:fcm-v} từ $U^{(t+1)}$
    \STATE \textbf{nếu} $\max |V^{(t+1)} - V^{(t)}| < \varepsilon$ \textbf{thì} dừng
  \ENDIF
\ENDFOR
\RETURN $U, V$ ở vòng cuối
\end{algorithmic}
\end{algorithm}
\subsection{Hàm rút gọn, hội tụ và vai trò của khởi tạo}
\label{sec:fcm-khoi-tao}

Với $V$ cố định, \eqref{eq:fcm-u} cho đúng ma trận $U$ làm $J_m$ nhỏ nhất. Thay~\eqref{eq:fcm-u} vào~\eqref{eq:fcm-j},
Hathaway và Bezdek~\cite{hathaway1995reformulation} đưa $J_m$ về một hàm chỉ của tâm:
\begin{equation}
  R_m(V) = \min_{U} J_m(U, V) = J_m\big(U^{*}(V), V\big)
  = \sum_{i=1}^{n} \left( \sum_{k=1}^{c} d_{ik}^{\,2/(1-m)} \right)^{1-m}.
  \label{eq:fcm-rut-gon}
\end{equation}
FCM làm giảm chính hàm này; Gröll và Jäkel~\cite{groll2005convergence} xem FCM trên $R_m$ là một phương pháp hạ dốc với bước
thay đổi. Mọi dãy lặp FCM hội tụ, ít nhất theo một dãy con, tới một cực tiểu địa phương hoặc một điểm yên ngựa của
$J_m$~\cite{bezdek1987counterexamples}; khởi tạo đủ gần một cực tiểu thì hội tụ về đúng cực tiểu đó, với tốc độ tuyến
tính~\cite{hathaway1988recent}. Từ đó suy ra điều kiện để một bộ khởi tạo tốt hơn có ích: \emph{$J_m$ phải có nhiều cực
tiểu, và các cách khởi tạo thông thường phải hay dừng ở cực tiểu kém}. Hathaway và Bezdek cũng ghi nhận trong thực tế FCM
khá ít phụ thuộc khởi tạo. Số cực tiểu của $J_m$ trên dữ liệu tín dụng vì vậy là câu hỏi thực nghiệm, được đo ở
mục~\ref{sec:vi-sao}; kết quả đo quyết định PSO nên được dùng vào việc gì (mục~\ref{sec:de-xuat}).

Mỗi vòng FCM tốn $O(ncD)$, nên một lần chạy tốn $O(I\,ncD)$ với $I$ là số vòng tới hội tụ. Tiểu luận cố định $m = 2$, giá
trị dùng phổ biến và cũng là giá trị của Qian và cộng sự~\cite{qian2022psofcm}; phần đối chiếu ở Chương~\ref{chap:thuc-nghiem}
kiểm tra lại với $m = 1{,}5$.

% ------------------------------------------------------------
\section{PSO gốc}
\label{sec:pso-goc}

Trong PSO~\cite{kennedy1995pso}, một bầy $N$ hạt cùng tìm cực tiểu của hàm $f$ trên $\mathbb{R}^{D}$. Hạt $j$ có vị trí
$\mathbf{x}_j = (x_{j1}, \dots, x_{jD})$ (nghiệm ứng viên), vị trí tốt nhất nó từng đi qua $\mathbf{p}_j$ (pbest) và vị trí tốt
nhất cả bầy từng gặp $\mathbf{g}$ (gbest). ``Vận tốc'' của hạt chỉ là bước dịch chuyển giữa 2 vòng liên tiếp,
$\Delta x_{jd}^{(t)} = x_{jd}^{(t)} - x_{jd}^{(t-1)}$. Ở vòng $t$, với mỗi chiều $d = 1, \dots, D$:
\begin{align}
  \Delta x_{jd}^{(t)} &= \omega(t)\, \Delta x_{jd}^{(t-1)}
     + c_1 r_{1jd} \big(p_{jd} - x_{jd}^{(t-1)}\big)
     + c_2 r_{2jd} \big(g_{d} - x_{jd}^{(t-1)}\big), \label{eq:pso-v}\\
  x_{jd}^{(t)} &= x_{jd}^{(t-1)} + \Delta x_{jd}^{(t)}, \label{eq:pso-x}
\end{align}
với $r_{1jd}, r_{2jd}$ là số ngẫu nhiên đều trên $[0, 1]$, rút mới ở mỗi vòng, mỗi hạt và mỗi chiều. Bước mới~\eqref{eq:pso-v}
gồm 3 phần: một phần $\omega(t)$ của bước cũ (quán tính), phần kéo về pbest (thành phần ``nhận thức'', hạt nhớ chỗ tốt của
mình) và phần kéo về gbest (thành phần ``xã hội'', hạt đi theo bầy).

Bản gốc năm 1995 không có hệ số $\omega$, tức $\omega(t) = 1$, và dùng $c_1 = c_2 = 2$~\cite{kennedy1995pso}. Vì vận tốc
cũ được giữ nguyên rồi cộng thêm, vận tốc có xu hướng tăng dần. Các cài đặt PSO thời kỳ đầu chặn vận tốc bằng một ngưỡng
tối đa $\Delta x_{\max}$ cho từng chiều (Shi và Eberhart dùng ngưỡng này trong thí nghiệm của mình~\cite{shi1998inertia}).
Tiểu luận đặt $\Delta x_{\max}$ bằng 0,2 lần độ rộng hộp tìm kiếm ở mỗi chiều. Không có $\omega < 1$ để giảm dần bước đi, bầy khó
thu hẹp vùng tìm kiếm về quanh gbest. Phiên bản cải tiến đầu tiên (mục~\ref{sec:pso-w}) nhắm vào điểm yếu này.

Thuật toán~\ref{alg:pso} là bản cài đặt dùng chung cho PSO gốc và 3 phiên bản có vận tốc (\texttt{pso\_co\_so.py});
các phiên bản chỉ khác nhau ở hàm $\omega(t)$. Trong thực nghiệm, PSO chạy đủ $T$ vòng, không dừng sớm, để ngân sách tính
toán là một tham số được kiểm soát. Chi phí là $N(T + 1)$ lần tính fitness $f$.

\begin{algorithm}[htbp]
\caption{PSO có vận tốc (tệp \texttt{pso\_co\_so.py}), dùng cho PSO gốc, PSO-W, PSO-S, PSO-LOG}
\label{alg:pso}
\begin{algorithmic}[1]
\REQUIRE Bầy ban đầu $\{\mathbf{x}_j\}_{j=1}^{N}$, hộp $[\mathbf{l}, \mathbf{u}]$, hàm fitness $f$ (cực tiểu hoá),
  lịch $\omega(t)$, số vòng $T$, $c_1 = c_2 = 2$
\ENSURE gbest $\mathbf{g}$ và $f(\mathbf{g})$
\STATE Kẹp $\mathbf{x}_j$ vào hộp; $\Delta\mathbf{x}_j \leftarrow \mathbf{0}$; $\mathbf{p}_j \leftarrow \mathbf{x}_j$;
  $\mathbf{g} \leftarrow \arg\min_j f(\mathbf{x}_j)$; $\Delta\mathbf{x}_{\max} \leftarrow 0{,}2(\mathbf{u} - \mathbf{l})$
\FOR{$t = 1$ \TO $T$}
  \FOR{mỗi hạt $j$}
    \STATE Cập nhật $\Delta\mathbf{x}_j$ theo~\eqref{eq:pso-v}, kẹp trong $\pm\Delta\mathbf{x}_{\max}$
    \STATE Cập nhật $\mathbf{x}_j$ theo~\eqref{eq:pso-x}, kẹp vào $[\mathbf{l}, \mathbf{u}]$
  \ENDFOR
  \STATE Chấm $f(\mathbf{x}_j)$ cho cả bầy; cập nhật $\mathbf{p}_j$ nếu tốt hơn
  \STATE $\mathbf{g} \leftarrow \arg\min_j f(\mathbf{p}_j)$ nếu tốt hơn gbest cũ
\ENDFOR
\RETURN $\mathbf{g}, f(\mathbf{g})$
\end{algorithmic}
\end{algorithm}

% ------------------------------------------------------------
\section{Các phiên bản cải tiến lấy từ các bài đã đọc}
\label{sec:pso-cai-tien}

Nhóm chọn 4 phiên bản theo 2 tiêu chí: là cải tiến được dùng rộng rãi (PSO-W, QPSO) hoặc đã được dùng với FCM trong
các bài nhóm đọc (PSO-S, PSO-LOG từ bài 02~\cite{qian2022psofcm}; QPSO từ bài 04~\cite{zhang2025qpsofcm}). Bảng~\ref{tab:ch2-phien-ban}
tóm tắt 5 thuật toán được so sánh.

\subsection{PSO-W: trọng số quán tính giảm tuyến tính}
\label{sec:pso-w}

Shi và Eberhart~\cite{shi1998inertia} thêm trọng số quán tính $\omega$ vào số hạng vận tốc cũ của~\eqref{eq:pso-v}. $\omega$
lớn giữ đà để bầy đi xa (khám phá), $\omega$ nhỏ làm hạt chậm lại quanh pbest và gbest (khai thác). Bài 1998 thử $\omega$
giảm tuyến tính trong khoảng rộng (từ 1,4 xuống 0). Lịch giảm tuyến tính từ $\omega_{\text{đầu}} = 0{,}9$ xuống
$\omega_{\text{cuối}} = 0{,}4$ được dùng rộng rãi về sau và thường được dẫn về nghiên cứu thực nghiệm năm 1999 của hai tác
giả~\cite{shi1999empirical}; nhóm không có toàn văn bài này nên chỉ dùng lịch theo cách các tài liệu về sau mô tả. Tiểu luận dùng lịch này:
\begin{equation}
  \omega(t) = \omega_{\text{đầu}} - \big(\omega_{\text{đầu}} - \omega_{\text{cuối}}\big)\,\frac{t - 1}{T - 1},
  \qquad t = 1, \dots, T.
  \label{eq:pso-omega}
\end{equation}
PSO gốc là trường hợp $\omega \equiv 1$. Lịch tuyến tính cũng là đối chứng L-PSO trong bài 02~\cite{qian2022psofcm} (công thức 30, với
$\omega_s = 1{,}2$, $\omega_e = 0{,}3$); tiểu luận giữ khoảng $0{,}9 \to 0{,}4$ cho PSO-W vì đây là giá trị chuẩn trong tài liệu,
còn PSO-S và PSO-LOG giữ nguyên tham số của bài 02.

\subsection{PSO-S và PSO-LOG: lịch quán tính phi tuyến}
\label{sec:pso-s-log}

Qian và cộng sự~\cite{qian2022psofcm} (bài 02) so 3 lịch $\omega$ cho PSO--FCM khi phân cụm tổ hợp quan sát của hệ thống
định vị BDS-3. Ngoài lịch tuyến tính, bài dùng 2 lịch phi tuyến, với $t = 1, \dots, T$:
\begin{align}
  \text{PSO-S (công thức 31):}\quad & \omega(t) = \omega_s - (\omega_s - \omega_e)\left(\frac{t}{T}\right)^{2}, \label{eq:omega-s}\\
  \text{PSO-LOG (công thức 32):}\quad & \omega(t) = \left[\ln\!\left(2{,}98 + \frac{t}{T}\right)\right]^{-z}, \label{eq:omega-log}
\end{align}
với $\omega_s = 1{,}2$, $\omega_e = 0{,}3$, $z = 1{,}51$ như bài; lịch logarit kế thừa từ ELPSO~\cite{zhang2020elpso}.
Dạng và hằng số ở~\eqref{eq:omega-log} theo bài 02; bài ELPSO gốc viết $\omega(t) = [\ln(2{,}1 + t)]^{-z}$ với $z = 1{,}05$
(công thức 12 trong~\cite{zhang2020elpso}), tức không chuẩn hoá $t$ theo $T$.
PSO-S giữ $\omega$ cao (trên 1 ở phần đầu) lâu hơn rồi giảm nhanh về $0{,}3$; PSO-LOG đi từ khoảng $0{,}87$ xuống khoảng
$0{,}61$, tức quán tính thay đổi ít hơn PSO-W. Bài không nêu $c_1, c_2$; tiểu luận dùng $c_1 = c_2 = 2$ như mọi phiên bản
có vận tốc khác để chỉ lịch $\omega$ thay đổi. Bài kết luận (tr.~21) S-PSO-FCM tốt nhất, nhưng theo Bảng~5 của chính bài,
ở 300 vòng LOG tốt nhất trên cả 3 chỉ số, còn S chỉ dẫn đầu ở 1.000 vòng. Thứ hạng đổi theo số vòng là một lý do để thực
nghiệm khảo sát số vòng.

\subsection{QPSO: PSO hành vi lượng tử}
\label{sec:qpso}

Trong QPSO~\cite{sun2004qpso}, hạt không có vận tốc mà được lấy mẫu quanh một điểm hút nằm giữa pbest và gbest, theo
phân phối mũ hai phía (phân phối Laplace). Ở dạng được dùng rộng rãi, mà Zhang và cộng sự dẫn về bài của Sun, Xu và Feng
cùng năm~\cite{sun2004global}, độ rộng của phân phối tỉ lệ với khoảng cách từ hạt tới trung bình các pbest
($\mathbf{m}_{\text{best}}$). Tiểu luận dùng dạng của
Zhang và cộng sự~\cite{zhang2025qpsofcm} (bài 04, QPSO--FCM). Ở mỗi vòng $t = 0, \dots, T - 1$, với mỗi hạt $j$ và mỗi
chiều $d = 1, \dots, D$:
\begin{align}
  \text{(công thức 16)}\quad & m_{\text{best},d} = \frac{1}{N} \sum_{j=1}^{N} p_{jd}, \label{eq:b-mbest}\\
  \text{(công thức 17, đã sửa)}\quad & \varphi_{jd} = \rho_{jd}\, p_{jd} + (1 - \rho_{jd})\, g_{d},
     \qquad \rho_{jd} \sim \mathcal{U}(0, 1), \label{eq:b-phi}\\
  \text{(công thức 18)}\quad & x_{jd} \leftarrow \varphi_{jd} \pm \gamma\, \big|m_{\text{best},d} - x_{jd}\big|\, \ln(1/q_{jd}), \label{eq:b-z}\\
  \text{(công thức 19)}\quad & \gamma = 0{,}9 \cdot \frac{T - t}{T} + 0{,}1. \label{eq:b-gamma}
\end{align}
Các phép nhân ở~\eqref{eq:b-phi} và~\eqref{eq:b-z} là phép nhân số thông thường trên từng toạ độ; $\rho_{jd}$ và $q_{jd}$ được
rút mới ở mỗi vòng, mỗi hạt và mỗi chiều, như $r_{1jd}, r_{2jd}$ ở~\eqref{eq:pso-v}. Dấu $\pm$ của từng toạ độ do một số đều
quyết định (cộng nếu số đó $\ge 0{,}5$). Với $q$ đều trên $(0, 1]$,
$\ln(1/q)$ có phân phối mũ với trung bình 1; ghép với dấu ngẫu nhiên, mỗi toạ độ mới có phân phối Laplace tâm
$\varphi_{jd}$, tham số tỉ lệ $\gamma\,|m_{\text{best},d} - x_{jd}|$. $\gamma$ giảm tuyến tính từ $1{,}0$ xuống
gần $0{,}1$ nên bước nhảy thu hẹp dần, tương tự $\omega$ giảm dần ở PSO-W. Mỗi pbest được cập nhật khi vị trí mới
tốt hơn. Bản in của công thức 17 là $\varphi_{ij} = \rho\, p_{ij} + (1 - \rho)\, p_{ij}$ (ký hiệu $\omega$ trong bài được đổi
thành $\rho$ để khỏi trùng với trọng số quán tính). Như vậy gbest không xuất hiện trong phương trình nào. Tiểu luận coi
đây là lỗi in và sửa theo QPSO chuẩn~\cite{sun2004qpso,sun2004global} như~\eqref{eq:b-phi}. Vị trí được kẹp vào hộp tìm
kiếm như các phiên bản khác (\texttt{thuat\_toan\_b.py}, hàm \texttt{qpso}).

\begin{table}[htbp]
\centering
\caption{Năm thuật toán PSO được so sánh}
\label{tab:ch2-phien-ban}
\footnotesize
\setlength{\tabcolsep}{4pt}
\begin{tabular}{@{}>{\raggedright\arraybackslash}p{2.1cm}>{\raggedright\arraybackslash}p{2.6cm}>{\raggedright\arraybackslash}p{4.4cm}>{\raggedright\arraybackslash}p{5.2cm}@{}}
\toprule
Tên & Nguồn & Cập nhật & Tham số \\
\midrule
PSO gốc & Kennedy, Eberhart 1995~\cite{kennedy1995pso} & \eqref{eq:pso-v}--\eqref{eq:pso-x} & $\omega = 1$, $c_1 = c_2 = 2$, $\Delta x_{\max} = 0{,}2 \times$ độ rộng hộp \\
PSO-W & Shi, Eberhart 1998, 1999~\cite{shi1998inertia,shi1999empirical} & như trên & $\omega$ theo~\eqref{eq:pso-omega}: $0{,}9 \to 0{,}4$ \\
PSO-S & Qian và cộng sự 2022~\cite{qian2022psofcm} & như trên & $\omega$ theo~\eqref{eq:omega-s}: $1{,}2 \to 0{,}3$ \\
PSO-LOG & Qian và cộng sự 2022~\cite{qian2022psofcm} & như trên & $\omega$ theo~\eqref{eq:omega-log}: $\approx 0{,}87 \to 0{,}61$ \\
QPSO & Sun và cộng sự 2004~\cite{sun2004qpso,sun2004global}; dạng của~\cite{zhang2025qpsofcm} & \eqref{eq:b-mbest}--\eqref{eq:b-gamma}, không có vận tốc & $\gamma$: $1{,}0 \to 0{,}1$ \\
\bottomrule
\end{tabular}
\par\smallskip\parbox{\linewidth}{\scriptsize Ghi chú: trong mỗi vai trò (tìm trọng số, mục~\ref{sec:pso-trong-so-pp};
tìm tâm, mục~\ref{sec:pso-khoi-tao}), mọi phiên bản dùng chung hộp tìm kiếm, cách khởi tạo bầy và hàm fitness.}
\end{table}


% ------------------------------------------------------------
\section{Các chỉ số đánh giá}
\label{sec:chi-so}

Các chỉ số trả lời 3 câu hỏi: phân hoạch tách rủi ro tốt tới đâu, phân hoạch tốt tới đâu về hình học, và tốn bao nhiêu
(Bảng~\ref{tab:ch2-chi-so}). Ba chỉ số chỉ dùng cho hướng tìm tâm được định nghĩa cùng hướng đó
(mục~\ref{sec:chi-so-tam}). Chỉ số phân cụm tính trên $(U, V)$ cuối của từng
lượt, trong không gian dữ liệu đã tiền xử lý (\texttt{danh\_gia.py}, \texttt{chi\_so\_danh\_gia.py}). Nhãn vỡ nợ
$y_i \in \{0, 1\}$ không đi vào thuật toán nào, chỉ dùng sau khi phân cụm. Nhãn cứng của điểm $i$ là
$\ell_i = \arg\max_k u_{ik}$, cụm cứng $k$ là $C_k = \{i : \ell_i = k\}$.

\subsection{Chỉ số tách rủi ro}
\label{sec:chi-so-rui-ro}

Phân cụm ở đây nhằm chia khách hàng thành các phân khúc có mức rủi ro khác nhau. Gọi $r_k$ là tỉ lệ vỡ nợ của cụm cứng
$C_k$. Hai chỉ số chính của tiểu luận đánh giá chính phân hoạch cứng:
\begin{itemize}
\item \emph{Chênh tỉ lệ vỡ nợ} $\Delta r = \max_k r_k - \min_k r_k$ trên các cụm không rỗng: phân khúc có ích cho quản
trị rủi ro khi các nhóm khác nhau rõ về tỉ lệ vỡ nợ.
\item \emph{AUC phân hoạch}: với điểm rủi ro $q_i = r_{\ell_i}$ (tỉ lệ vỡ nợ của cụm chứa $i$),
\begin{equation}
  \mathrm{AUC} = \frac{1}{n_1 n_0} \sum_{i : y_i = 1} \; \sum_{j : y_j = 0}
     \Big( \mathbb{1}[q_i > q_j] + \tfrac{1}{2}\,\mathbb{1}[q_i = q_j] \Big),
  \label{eq:chi-so-auc}
\end{equation}
với $n_1, n_0$ là số khách hàng vỡ nợ và không vỡ nợ~\cite{fawcett2006roc}: xác suất một khách hàng vỡ nợ chọn ngẫu nhiên
nằm ở phân khúc rủi ro cao hơn một khách hàng không vỡ nợ chọn ngẫu nhiên. Vì $q_i$ tăng theo tỉ lệ vỡ nợ của cụm, chỉ số này
luôn từ $0{,}5$ trở lên, bằng $0{,}5$ khi mọi cụm có cùng tỉ lệ vỡ nợ.
\end{itemize}
\subsection{Chỉ số về kết quả phân cụm}
Hệ số phân hoạch (PC) và entropy phân hoạch chuẩn hoá (PE)~\cite{bezdek1981pattern} đo độ ``rõ'' của phân hoạch:
\begin{equation}
  \mathrm{PC} = \frac{1}{n} \sum_{i=1}^{n} \sum_{k=1}^{c} u_{ik}^{2} \in \left[\tfrac{1}{c}, 1\right], \qquad
  \mathrm{PE} = -\frac{1}{n \log_2 c} \sum_{i=1}^{n} \sum_{k=1}^{c} u_{ik} \log_2 u_{ik} \in [0, 1].
  \label{eq:chi-so-pc-pe}
\end{equation}
PC lớn, PE nhỏ là phân hoạch gần cứng; PC $= 1/c$ và PE $= 1$ khi mọi độ thuộc bằng $1/c$, dấu hiệu của nghiệm suy biến có
các tâm trùng nhau. Tiểu luận phân biệt 2 mức: nghiệm \emph{sụp} khi mọi tâm dồn về trung bình dữ liệu (độ lớn tâm dưới
$10^{-3}$ trên dữ liệu z-score), và nghiệm \emph{gần suy biến} khi chỉ một số tâm gần trùng nhau (khoảng cách nhỏ nhất giữa
2 tâm rất nhỏ so với độ trải của dữ liệu). Chỉ số Xie--Beni (XB)~\cite{xie1991validity} đo tỉ số giữa độ nén và độ tách:
\begin{equation}
  \mathrm{XB} = \frac{\sum_{i=1}^{n} \sum_{k=1}^{c} u_{ik}^{m} \lVert \mathbf{x}_i - \mathbf{v}_k \rVert^{2}}
                     {n \cdot \min_{k \ne l} \lVert \mathbf{v}_k - \mathbf{v}_l \rVert^{2}}.
  \label{eq:chi-so-xb}
\end{equation}
XB nhỏ là tốt. Mẫu số là bình phương khoảng cách nhỏ nhất giữa 2 tâm, nên khi 2 tâm gần trùng, một thay đổi rất nhỏ
của tâm làm XB thay đổi rất lớn (mã nguồn trả vô cùng khi trùng hẳn). PC và XB là 2 chỉ số bài 04 dùng (công thức 21--22
trong~\cite{zhang2025qpsofcm}). Hệ số Silhouette~\cite{rousseeuw1987silhouettes} đánh giá phân hoạch cứng. Với $a(i)$ là
khoảng cách trung bình từ $i$ tới các điểm cùng cụm và $b(i)$ là khoảng cách trung bình nhỏ nhất tới các điểm của một cụm
khác,
\begin{equation}
  s(i) = \frac{b(i) - a(i)}{\max\{a(i), b(i)\}}, \qquad \mathrm{Sil} = \frac{1}{|S|} \sum_{i \in S} s(i) \in [-1, 1],
  \label{eq:chi-so-sil}
\end{equation}
với $S$ là một mẫu cố định $3.000$ điểm (toàn bộ nếu $n \le 3.000$, seed 0) dùng chung cho mọi phương pháp và mọi lượt
trên cùng bộ dữ liệu.
Silhouette cũng là fitness của đề xuất (mục~\ref{sec:pso-trong-so-pp}, mẫu 1.000 điểm khác mẫu đánh giá), nên
khi so PSO với đối chứng, Silhouette chỉ là chỉ số phụ; 2 chỉ số tách rủi ro không nằm trong fitness.
\subsection{Chỉ số về chi phí}

Thước đo chi phí là \emph{thời gian thực} (\texttt{perf\_counter}) của cả quá trình, gồm khởi tạo, tìm kiếm (nếu có) và FCM,
trên cùng máy và cùng thư viện (mục~\ref{sec:thiet-lap}).

\begin{table}[htbp]
\centering
\caption{Các chỉ số đánh giá}
\label{tab:ch2-chi-so}
\footnotesize
\setlength{\tabcolsep}{4pt}
\begin{tabular}{@{}>{\raggedright\arraybackslash}p{3.8cm}>{\raggedright\arraybackslash}p{1.6cm}>{\raggedright\arraybackslash}p{1.4cm}>{\raggedright\arraybackslash}p{7.6cm}@{}}
\toprule
Chỉ số & Tốt khi & Dùng nhãn & Vai trò \\
\midrule
Chênh tỉ lệ vỡ nợ $\Delta r$ & lớn & có & \textbf{chính}: các phân khúc khác nhau về rủi ro tới mức nào \\
AUC phân hoạch \eqref{eq:chi-so-auc} & lớn & có & \textbf{chính}: khả năng xếp hạng rủi ro của phân hoạch cứng; $\ge 0{,}5$ \\
Silhouette \eqref{eq:chi-so-sil} & lớn & không & chất lượng phân hoạch cứng; là fitness khi tìm trọng số \\
PC, PE \eqref{eq:chi-so-pc-pe} & PC lớn, PE nhỏ & không & độ rõ của phân hoạch mờ; phát hiện nghiệm suy biến \\
XB \eqref{eq:chi-so-xb} & nhỏ & không & nén/tách; không ổn định khi 2 tâm gần trùng \\
Thời gian thực & nhỏ & không & chi phí \\
\bottomrule
\end{tabular}
\end{table}


% ------------------------------------------------------------
\section{PSO tìm tâm khởi tạo cho FCM}
\label{sec:pso-khoi-tao}

Cách ghép PSO với FCM phổ biến trong tài liệu (mục~\ref{sec:metaheuristic}) là dùng PSO tìm tâm, và đây là hướng nhóm thử
đầu tiên. Nhóm dùng dạng đơn giản nhất của cách này: mỗi hạt là một bộ tâm $V \in \mathbb{R}^{c \times D}$, fitness là
$R_m(V)$ theo~\eqref{eq:fcm-rut-gon} (cách mã hoá tương tự Runkler và Katz~\cite{runkler2006psofcm}), mỗi hạt ban đầu là $c$
điểm dữ liệu khác nhau chọn ngẫu nhiên, hộp tìm kiếm là $[\min_i x_{id}, \max_i x_{id}]$ theo từng thuộc tính. PSO chạy đủ
$T$ vòng, rồi gbest làm tâm khởi tạo cho FCM chạy tới hội tụ (Thuật toán~\ref{alg:pso-khoi-tao}, \texttt{pso\_khoi\_tao.py});
cách ghép tuần tự này là cách bài 04 dùng cho QPSO--FCM~\cite{zhang2025qpsofcm}. Năm phiên bản được khảo sát với số hạt
$N \in \{10, 20, 30\}$ và số vòng $T \in \{20, 30, 50\}$.

\begin{algorithm}[htbp]
\caption{PSO tìm tâm khởi tạo cho FCM (\texttt{pso\_khoi\_tao.py})}
\label{alg:pso-khoi-tao}
\begin{algorithmic}[1]
\REQUIRE Dữ liệu $X$, số cụm $c$, $m = 2$, phiên bản PSO, số hạt $N$, số vòng $T$
\ENSURE Tâm $V^{*}$, độ thuộc $U^{*}$
\STATE Mỗi hạt $j = 1, \dots, N$: chọn $c$ điểm dữ liệu khác nhau ngẫu nhiên làm $\mathbf{x}_j$
\STATE Chạy phiên bản PSO đã chọn đủ $T$ vòng với fitness $R_m$; được gbest $\mathbf{g}$
\STATE $(U^{*}, V^{*}) \leftarrow$ FCM (Thuật toán~\ref{alg:fcm}) khởi động từ tâm $\mathbf{g}$, chạy tới hội tụ
\RETURN $U^{*}, V^{*}$
\end{algorithmic}
\end{algorithm}

Vì PSO và FCM làm giảm cùng một hàm $R_m$, PSO chỉ có thể giúp khi $R_m$ có nhiều cực tiểu (mục~\ref{sec:fcm-khoi-tao}).
Đối chứng là 3 cách khởi tạo không tốn tìm kiếm, cùng kết thúc bằng FCM: độ thuộc $U$ ngẫu nhiên (dạng gốc của Bezdek),
$c$ điểm dữ liệu ngẫu nhiên (đúng một hạt PSO trước khi bay, nên khác biệt chỉ là phần tìm kiếm) và
k-means++~\cite{arthur2007kmeanspp} (tâm sau chọn với xác suất tỉ lệ bình phương khoảng cách tới tâm gần nhất đã chọn).

\emph{Chỉ số riêng của hướng này.} Hai chỉ số đầu đo chất lượng tối ưu. Gọi $J^{\star}$ là $J_m$ nhỏ nhất trên mọi lượt của mọi phương pháp
ở cùng cấu hình (bộ dữ liệu, $m$, $c$). \emph{Độ lệch $J$ sau PSO} là $100\,(R_m(\mathbf{g}) - J^{\star}) / J^{\star}$ (\%), đo
tại gbest trước khi FCM chạy, so các phiên bản PSO như những thuật toán tối ưu. \emph{Tỉ lệ đạt $J$ tốt nhất} đo kết quả sau
FCM: một lượt đạt nếu $J_m \le (1 + 10^{-4})\,J^{\star}$; ngưỡng $10^{-4}$ (tương đối) cũng là ngưỡng để coi 2 lượt cho
``cùng nghiệm''. Chỉ số thứ ba, theo cách các bài đã đọc đánh giá phân cụm, là \emph{AUC tách rủi ro}: \eqref{eq:chi-so-auc} với điểm rủi ro là độ thuộc vào
cụm có tỉ lệ vỡ nợ cao nhất, $q_i = u_{ik^{*}}$, $k^{*} = \arg\max_k r_k$. Độ thuộc đo độ gần tâm và không nhất thiết tăng
theo rủi ro, nên chỉ số này có thể nhỏ hơn $0{,}5$ khi các tâm gần trùng.
\label{sec:chi-so-tam}

% ------------------------------------------------------------
\section{Trọng số thuộc tính trong FCM}
\label{sec:trong-so-fcm}

Khi PSO tìm tâm không cải thiện FCM (mục~\ref{sec:chuyen-huong}), nhóm xét điểm yếu thứ hai của FCM nêu ở
mục~\ref{sec:phan-cum-mo}: mọi thuộc tính được coi ngang nhau. Mục này chỉ ra vì sao FCM không tự sửa được điểm yếu này, và
cách sửa bằng công thức có giới hạn gì; đó là chỗ đề xuất ở mục~\ref{sec:de-xuat} nhắm vào.

\subsection{Khoảng cách có trọng số thuộc tính}
\label{sec:kc-trong-so}

FCM ở~\eqref{eq:fcm-j} dùng khoảng cách Euclid trên dữ liệu đã chuẩn hoá z-score, tức mọi thuộc tính có ảnh hưởng ngang
nhau tới khoảng cách. Với vector trọng số không âm $\symbf{\lambda} = (\lambda_1, \dots, \lambda_D)$, khoảng cách có trọng
số là
\begin{equation}
  d_{ik}^{2}(\symbf{\lambda}) = \sum_{j=1}^{D} \frac{\lambda_j}{\bar{\lambda}} \, (x_{ij} - v_{kj})^{2},
  \qquad \bar{\lambda} = \frac{1}{D} \sum_{j=1}^{D} \lambda_j .
  \label{eq:kc-trong-so}
\end{equation}
Chia cho $\bar{\lambda}$ làm khoảng cách chỉ phụ thuộc tỉ lệ giữa các trọng số, và với $\symbf{\lambda} = \mathbf{1}$ ta có
lại FCM thường. Với $\symbf{\lambda}$ cố định, FCM theo~\eqref{eq:kc-trong-so} chính là FCM thường chạy trên dữ liệu đã nhân
từng cột với $s_j = \sqrt{\lambda_j / \bar{\lambda}}$, nên mọi tính chất ở mục này giữ nguyên; tâm báo cáo được tính lại
trên dữ liệu gốc bằng~\eqref{eq:fcm-v} với độ thuộc cuối (hàm \texttt{phan\_hoach\_trong\_so}).

FCM không tự chọn được $\symbf{\lambda}$. Nếu cực tiểu hoá cả $J_m$ theo $\symbf{\lambda}$ với tổng trọng số cố định,
$J_m$ tuyến tính theo $\symbf{\lambda}$ nên cực tiểu nằm ở một đỉnh: dồn toàn bộ trọng số vào thuộc tính có độ phân tán
trong cụm nhỏ nhất, và phân cụm chỉ còn theo một thuộc tính. Vì vậy cần hoặc một số mũ để làm trơn bài toán
(mục~\ref{sec:wfcm}), hoặc một tiêu chí khác $J_m$ để chọn trọng số (mục~\ref{sec:pso-trong-so-pp}).

\subsection{Trọng số tính bằng công thức (WFCM) và giới hạn}
\label{sec:wfcm}

Keller và Klawonn~\cite{keller2000weighting} làm trơn bài toán ở mục~\ref{sec:kc-trong-so} bằng một số mũ $\beta > 1$: khoảng
cách dùng $w_j^{\beta}$ thay cho $\lambda_j$, với $\sum_j w_j = 1$. Khi đó cực tiểu theo $\mathbf{w}$ (giữ $U, V$) có dạng
đóng
\begin{equation}
  w_j \propto \left( \sum_{i=1}^{n} \sum_{k=1}^{c} u_{ik}^{m} (x_{ij} - v_{kj})^{2} \right)^{-1/(\beta - 1)},
  \label{eq:wfcm-w}
\end{equation}
và thuật toán WFCM luân phiên cập nhật $U$, $V$ và $\mathbf{w}$. Tiểu luận dùng $\beta = 2$, khởi tạo $w_j = 1/D$ và tâm
k-means++, dừng khi tâm thay đổi dưới $10^{-5}$ (\texttt{wfcm\_giai\_tich} trong \texttt{cach\_ket\_hop.py}); trọng số
tương đương theo~\eqref{eq:kc-trong-so} là $\lambda_j = w_j^{2}$. Bài 06~\cite{liu2021msgwo} dùng một biến thể có trọng số
riêng cho từng cụm, cũng tính bằng công thức, còn tâm do thuật toán bầy sói xám cải tiến (MSGWO) tìm.

Công thức~\eqref{eq:wfcm-w} cho thuộc tính nào làm các cụm chặt hơn thì nặng hơn, tức vẫn theo tiêu chí ``cụm chặt'' của
$J_m$ và không xét các cụm có tách nhau hay không. Trên dữ liệu có nhóm thuộc tính tương quan cao và ít biến động (chẳng hạn dư nợ các
tháng của bộ Taiwan Default), công thức có thể dồn trọng số vào nhóm đó dù nhóm này không tạo ra cấu trúc nhóm có ý nghĩa.
WFCM giải tích rẻ, chỉ tốn cỡ một lần chạy FCM, nên là đối chứng cần thiết: PSO chỉ đáng dùng khi nó tốt hơn công thức này.

% ------------------------------------------------------------
\section{Đề xuất cải tiến: PSO tìm trọng số thuộc tính cho FCM}
\label{sec:pso-trong-so-pp}
\label{sec:de-xuat}

\subsection{Ý tưởng}

Trên dữ liệu tín dụng, nhiều chỉ tiêu đo gần cùng một thứ (các tỉ số nợ trên vốn chủ sở hữu và trên tổng tài sản; dư nợ
các tháng liên tiếp), còn một số chỉ tiêu gần như không liên quan tới cấu trúc nhóm, nên khi trọng số bằng nhau, các thuộc tính
thừa có thể chi phối phân hoạch. Trọng số là một vector $D$ chiều ($D$ từ 14 tới 94 trên các bộ đã thử), quá
nhiều để quét lưới; chất lượng phân hoạch theo trọng số không có đạo hàm vì mỗi lần đánh giá là một lần chạy FCM; và công
thức giải tích chỉ tối ưu độ chặt (mục~\ref{sec:wfcm}). Ở vai trò này, một metaheuristic như PSO tối ưu thứ FCM không tự
tối ưu được; khi tìm tâm, PSO tối ưu đúng hàm $R_m$ mà FCM đã tự làm giảm (mục~\ref{sec:pso-khoi-tao}). Nhóm chuyển sang
hướng này sau khi PSO tìm tâm không cải thiện FCM trên dữ liệu tín dụng (mục~\ref{sec:chuyen-huong}).

Trong 10 bài đã đọc, đề xuất dựa trên 3 bài. Bài 06~\cite{liu2021msgwo} nêu việc coi mọi thuộc tính ngang nhau là một điểm
yếu của FCM và cho thấy đánh trọng số đóng góp nhiều hơn tìm tâm (mục~\ref{sec:tinh-hinh-nghien-cuu}), nhưng trọng số ở đó tính
bằng công thức còn bầy sói tìm tâm; ở đây tâm khởi tạo do k-means++ chọn và PSO chỉ tìm trọng số. Bài 02 và 04 cung cấp
các phiên bản PSO (mục~\ref{sec:pso-cai-tien}), được lắp vào khung tìm trọng số thay vì khung tìm tâm. Ngoài 10 bài, Lu
và cộng sự~\cite{lu2011psovw} và O'Neill và cộng sự~\cite{oneill2018psofw} đã dùng PSO tìm trọng số thuộc tính cho phân cụm
cứng; tiểu luận đưa cách làm đó vào FCM trên dữ liệu tín dụng và so 5 phiên bản PSO trong vai trò này. Cách làm của nhóm khác O'Neill và cộng sự ở chỗ bài đó cho PSO
vừa chọn vừa đánh trọng số thuộc tính, tính Silhouette trên chính tập thuộc tính đã đánh trọng số (chỉ thành phần connectedness dùng toàn bộ
thuộc tính gốc, mục III.B), và ghi nhận ở phần kết luận rằng một số nghiệm sinh ra cụm tầm thường, vấn đề bài để ngỏ.

\subsection{Mã hoá và fitness}

Mỗi hạt là vector trọng số $\symbf{\lambda} \in [0{,}05;\ 1]^D$, dùng trong khoảng cách~\eqref{eq:kc-trong-so}. Cận dưới
$0{,}05$ giữ mọi thuộc tính trong bài toán, vì mục tiêu là đánh trọng số chứ không chọn bỏ thuộc tính. Fitness của một hạt tính
theo 3 bước (\texttt{bai\_toan\_trong\_so} trong \texttt{cach\_ket\_hop.py}):
\begin{enumerate}[leftmargin=*]
  \item chạy FCM (tối đa 200 vòng) với khoảng cách~\eqref{eq:kc-trong-so}, từ một bộ tâm k-means++ $V^{(0)}$ cố định cho
  cả lượt;
  \item lấy phân hoạch cứng, tính Silhouette~\eqref{eq:chi-so-sil} \emph{trên không gian gốc} với một mẫu cố định 1.000 điểm;
  \item $f(\symbf{\lambda}) = -\mathrm{Sil} + P$, với phạt $P = \max(0,\ 0{,}05 - n_{\min}/n) / 0{,}05$ khi cụm cứng nhỏ nhất
  có ít hơn 5\% số khách hàng.
\end{enumerate}
Ba lựa chọn thiết kế trong bài toán này là phần cải tiến của nhóm:
\begin{enumerate}[label=(\alph*), leftmargin=*]
\item \emph{Silhouette tính trên không gian gốc.} Silhouette đo độ tách giữa các cụm, bù cho tiêu chí độ chặt của $J_m$. Nếu
tính trên không gian đã nhân trọng số, PSO có thể tăng chỉ số chỉ bằng cách thu nhỏ các chiều khó tách; tính trên không gian
gốc buộc phân hoạch phải tốt đối với dữ liệu thật.
\item \emph{Phạt cụm dưới 5\%.} Silhouette có thể cao khi một cụm chỉ gồm vài điểm ngoại lai; phân khúc như vậy không dùng
được cho chính sách tín dụng. Đây là cách nhóm xử lý vấn đề cụm tầm thường mà O'Neill và cộng sự để ngỏ.
\item \emph{Hạt khởi đầu $\symbf{\lambda} = \mathbf{1}$ và tâm khởi tạo cố định.} Nghiệm của PSO không bao giờ kém FCM thường
theo fitness, và mọi khác biệt giữa các trọng số chỉ do trọng số, không do khởi tạo.
\end{enumerate}

\subsection{Thuật toán và tham số}

Hạt đầu tiên của bầy là $\symbf{\lambda} = \mathbf{1}$ như lựa chọn (c), $N - 1$ hạt còn lại rút đều trong hộp. Bất kỳ phiên bản nào ở mục~\ref{sec:pso-cai-tien} cũng lắp được
vào khung này mà không đổi gì khác: 4 phiên bản có vận tốc chạy Thuật toán~\ref{alg:pso} với lịch $\omega(t)$ của mình,
$c_1 = c_2 = 2$; QPSO chạy~\eqref{eq:b-mbest}--\eqref{eq:b-gamma}. Mọi phiên bản chạy đủ 49 vòng sau khởi tạo, không dừng
sớm, nên ngân sách là $50N$ lần tính fitness. Tiểu luận dùng 2 ngân sách: 500 lần (10 hạt) và 1.000 lần (20 hạt). Sau PSO,
FCM chạy tới hội tụ với trọng số gbest từ cùng $V^{(0)}$ (Thuật toán~\ref{alg:pso-trong-so}, \texttt{pso\_trong\_so.py}).
Thí nghiệm chính so với các đối chứng dùng PSO-W với 500 lần; thí nghiệm so phiên bản dùng cả 5 phiên bản ở 2 ngân sách.

\begin{algorithm}[htbp]
\caption{PSO tìm trọng số thuộc tính cho FCM (\texttt{pso\_trong\_so.py})}
\label{alg:pso-trong-so}
\begin{algorithmic}[1]
\REQUIRE $X \in \mathbb{R}^{n \times D}$ đã chuẩn hoá, số cụm $c$, $m = 2$, phiên bản PSO, số hạt $N$, ngân sách $50N$ lần tính fitness
\ENSURE Trọng số $\symbf{\lambda}^{*}$ và phân hoạch $(U, V)$
\STATE Chọn $V^{(0)}$ bằng k-means++; chọn mẫu Silhouette $S$ gồm 1.000 điểm
\STATE Khởi tạo bầy: hạt 1 là $\mathbf{1}$, $N - 1$ hạt rút đều trong $[0{,}05;\ 1]^D$; tính $f$ cho mọi hạt
\FOR{$t = 1$ \TO $49$}
  \STATE Di chuyển bầy theo phiên bản đã chọn (Thuật toán~\ref{alg:pso} hoặc QPSO); kẹp vào hộp
  \STATE Với mỗi hạt: FCM với khoảng cách~\eqref{eq:kc-trong-so} từ $V^{(0)}$; $f = -\mathrm{Sil}_S + P$
  \STATE Cập nhật pbest, gbest
\ENDFOR
\STATE $\symbf{\lambda}^{*} \leftarrow$ gbest; FCM tới hội tụ với $\symbf{\lambda}^{*}$ từ $V^{(0)}$
\RETURN $\symbf{\lambda}^{*}$, $U$ cuối và $V$ tính trên không gian gốc
\end{algorithmic}
\end{algorithm}

Mỗi lần tính fitness gồm một lần chạy FCM (tối đa 200 vòng, $O(ncD)$ mỗi vòng) và một lần tính Silhouette trên 1.000 điểm
($O(10^6)$ phép tính khoảng cách cho mỗi phân hoạch), nên một lượt PSO tốn cỡ vài trăm lần chạy FCM. Đây là chi phí trả một
lần khi dựng phân khúc.

\subsection{Các đối chứng}

Ba đối chứng tách 3 nguồn lợi ích có thể có:
\begin{itemize}
\item \emph{FCM k-means++} ($\symbf{\lambda} = \mathbf{1}$, cùng $V^{(0)}$) cho biết việc đánh trọng số có ích không;
\item \emph{WFCM giải tích} (mục~\ref{sec:wfcm}) cho biết có cần tìm kiếm không, hay một công thức rẻ là đủ;
\item \emph{Tìm ngẫu nhiên} rút $50N - 1$ vector trong cùng hộp, cộng $\mathbf{1}$, giữ vector có fitness nhỏ nhất, với cùng
$V^{(0)}$, cùng mẫu Silhouette và cùng ngân sách như PSO khi cùng seed (\texttt{ngau\_nhien\_trong\_so}). Phần chênh giữa PSO
và tìm ngẫu nhiên là phần do cơ chế bầy đàn đem lại, không do việc đánh trọng số hay do ngân sách tính toán.
\end{itemize}
Để kiểm tra lợi ích có giữ ngoài dữ liệu đã dùng để tìm trọng số, trọng số học trên một mẫu còn được áp lên một mẫu độc lập
(mục~\ref{sec:kiem-chung}).

\subsection{Quy trình áp dụng}
\label{sec:quy-trinh}

Đề xuất được dùng trong một quy trình 3 bước (Thuật toán~\ref{alg:chan-doan}), rút ra từ chính 2 giai đoạn của tiểu luận:
trước khi nghĩ tới PSO cho tâm, đo xem $J_m$ có nhiều cực tiểu không; dùng PSO cho trọng số; và chỉ áp trọng số khi nó cải
thiện phân khúc trên dữ liệu có nhãn lịch sử không dùng để tìm trọng số.

\begin{algorithm}[htbp]
\caption{Quy trình áp dụng đề xuất khi phân cụm khách hàng bằng FCM}
\label{alg:chan-doan}
\begin{algorithmic}[1]
\REQUIRE Dữ liệu $X$ đã tiền xử lý, số cụm $c$, số mũ $m$, số lượt $R$ (ví dụ $R = 30$), ngưỡng gộp $\tau = 10^{-4}$; nhãn
  lịch sử $y$ (chỉ dùng ở bước 3)
\ENSURE Cách khởi tạo, trọng số thuộc tính và phân hoạch
\STATE \textbf{Bước 1, chẩn đoán số cực tiểu:} chạy FCM từ $R$ khởi tạo $U$ ngẫu nhiên và $R$ khởi tạo k-means++; sắp $2R$
  giá trị $J_m$ cuối tăng dần, gộp các giá trị lệch nhau dưới $(1 + \tau)$ lần vào cùng mức, được $L_1 < \dots < L_K$
\IF{các lượt ở mức $L_1$ là nghiệm sụp (mọi tâm gần trung bình dữ liệu, $\mathrm{PC} \approx 1/c$)}
  \STATE Đổi tiền xử lý, $m$ hoặc $c$; khởi tạo không cứu được
\ELSIF{$K = 1$ hoặc k-means++ hầu như luôn đạt $L_1$}
  \STATE Khởi tạo bằng k-means++; \textbf{không} dùng PSO tìm tâm
\ELSE
  \STATE Dùng k-means++ nhiều lần, giữ lượt có $J_m$ nhỏ nhất; chỉ cân nhắc PSO tìm tâm nếu cách này vẫn hay trượt $L_1$
\ENDIF
\STATE \textbf{Bước 2, tìm trọng số:} chạy PSO tìm trọng số (Thuật toán~\ref{alg:pso-trong-so}) với tâm khởi tạo đã chọn;
  so fitness với WFCM giải tích và với tìm ngẫu nhiên cùng ngân sách
\STATE \textbf{Bước 3, kiểm bằng nhãn lịch sử:} so $\Delta r$ và AUC phân hoạch của phân hoạch có trọng số với FCM thường,
  trên một mẫu không dùng để tìm trọng số; chỉ dùng trọng số nếu phân khúc tách rủi ro tốt hơn
\RETURN cách khởi tạo, $\symbf{\lambda}^{*}$ (hoặc $\mathbf{1}$) và phân hoạch tương ứng
\end{algorithmic}
\end{algorithm}

\emph{Bước 1} dùng 2 kiểu khởi tạo có vai trò khác nhau. Khởi tạo $U$ ngẫu nhiên là dạng gốc của FCM, tâm ban đầu nằm gần
trung bình dữ liệu, nên có thể lộ ra các mức mà k-means++ không tới, kể cả nghiệm sụp. Khởi tạo k-means++ là cách sẽ dùng nếu
không cần PSO, nên tần suất nó đạt $L_1$ trả lời trực tiếp câu hỏi có cần bộ khởi tạo tốt hơn không. Số mức $K$ là số cực
tiểu \emph{quan sát được}, không phải số cực tiểu thật; ngưỡng $\tau$ trùng ngưỡng của tỉ lệ đạt $J$ tốt nhất, và cách gộp
giống hàm \texttt{so\_cuc\_tieu} trong \texttt{phan\_tich.py}. Bước này tốn $2R$ lần chạy FCM, chỉ làm một lần cho mỗi bộ
$(X, m, c)$, và tự cho luôn một nghiệm dùng được. \emph{Bước 2} dùng 2 đối chứng để biết có cần PSO không: nếu WFCM giải
tích đạt cùng fitness thì dùng công thức; nếu tìm ngẫu nhiên đạt cùng fitness thì phiên bản PSO không đóng góp gì.
\emph{Bước 3} cần thiết vì fitness không nhìn nhãn: một phân hoạch đẹp hơn về hình học không nhất thiết tách rủi ro tốt hơn.
Tổ chức tín dụng thường có nhãn vỡ nợ lịch sử; dùng nhãn này để \emph{chọn} có áp trọng số hay không, trên một mẫu độc lập
với mẫu tìm trọng số, giữ cho bước phân cụm vẫn không giám sát.

Đóng góp của đề xuất vì vậy gồm 2 phần: bài toán tìm trọng số với 3 lựa chọn thiết kế cho fitness, và quy trình quyết định
PSO dùng vào việc gì. Đề xuất không tạo ra phiên bản PSO mới; mọi phiên bản ở mục~\ref{sec:pso-cai-tien} đều lắp được vào nó.
Chương~\ref{chap:thuc-nghiem} đánh giá đề xuất: bước 1 ở mục~\ref{sec:diem-xuat-phat}, bước 2 ở mục~\ref{sec:pso-trong-so}
và~\ref{sec:phien-ban-trong-so}, bước 3 ở mục~\ref{sec:kiem-chung}.

\section*{Tiểu kết Chương 2}
\addcontentsline{toc}{section}{Tiểu kết Chương 2}

Chương trình bày FCM, PSO gốc, 4 phiên bản cải tiến (PSO-W; PSO-S và PSO-LOG từ bài 02; QPSO từ bài 04) và các chỉ số
đánh giá, rồi 2 cách dùng PSO theo thứ tự nhóm thực hiện. Hai chỉ số chính là chênh tỉ lệ vỡ nợ và AUC phân hoạch, đều dùng
nhãn và không nằm trong fitness nào. Ở hướng ban đầu, mỗi hạt là một bộ tâm và fitness là hàm rút gọn $R_m$; vì PSO và FCM làm giảm cùng một hàm, PSO chỉ
có thể giúp khi $R_m$ có nhiều cực tiểu. Hướng thứ hai nhắm vào điểm yếu FCM không tự sửa được: cực tiểu hoá $J_m$ theo trọng
số cho nghiệm suy biến, còn công thức của WFCM chỉ tối ưu độ chặt của cụm. Đề xuất cải tiến của nhóm dùng PSO tìm trọng số
thuộc tính, với 3 lựa chọn thiết kế cho fitness (Silhouette trên không gian gốc, phạt cụm dưới 5\%, hạt khởi đầu là FCM
thường), lắp được mọi phiên bản PSO, ở 2 ngân sách 500 và 1.000 lần tính fitness, và được dùng trong quy trình 3 bước:
chẩn đoán số cực tiểu, tìm trọng số, kiểm bằng nhãn lịch sử trên mẫu độc lập. Ba đối chứng tách 3 nguồn lợi ích: FCM thường,
WFCM giải tích và tìm ngẫu nhiên cùng ngân sách.
